\documentclass[11pt,a4paper]{article}
\usepackage[margin=25mm]{geometry}
\usepackage{amsmath,amssymb,amsthm,hyperref}
\newtheorem{theorem}{Theorem}\newtheorem{lemma}[theorem]{Lemma}\newtheorem{proposition}[theorem]{Proposition}
\newcommand{\gen}[1]{\langle#1\rangle}\newcommand{\F}{\mathbb F}\newcommand{\Aut}{\operatorname{Aut}}
\title{Flag completion and Wilson symmetry in the four-type construction}
\author{Human Compression project: L023 independent Section 5 recovery}\date{}
\begin{document}\maketitle

This paper concerns Construction 6, Lemmas 7--8 and Theorem 10 of
Leemans--Stokes, arXiv:2208.08215v2, physical pages 10--12. The source
identifies two notions of map: a faithful vertex--edge--face incidence
map and an arbitrary quotient of the universal map group. They need not
retain the same incidence information. Our counterexamples below are
actual finite, faithful, thin incidence maps coming from surface embeddings;
they therefore do not depend on that ambiguity. The conclusion about
trialities without dualities survives for their four-type pregeometries,
but the asserted universal flag completion does not.

\section{Local chambers and the precise reconstruction hypothesis}

Write $a,b,c$ for a marked triple of nonidentity involutions, with
$ac=ca$ and $V=\gen{a,c}$ of order four. Set
\[
 H_0=\gen{a,b},\quad H_1=V,\quad H_2=\gen{b,c},
 \quad H_3=\gen{b,ac}.
\]
Objects of type $i$ are left cosets of $H_i$; incidence means nonempty
intersection. Types $0,1,2,3$ are respectively faces, edges, vertices
and Petrie polygons. We distinguish a full-type chamber from a maximal
flag which might fail to have every type.

\begin{lemma}\label{local}
Suppose $G=\gen{a,b,c}$ and each of $H_0,H_2,H_3$ is
proper. Every flag containing an edge extends to a full-type chamber.
There are exactly two $G$-orbits of such chambers, both free; if $G$
is finite their total size is $2|G|$. The
pregeometry is a geometry if and only if every mutually incident
face--vertex--Petrie triple has a common incident edge.
\end{lemma}
\begin{proof}
If $H_0$ contained $c$ it would be all of $G$, so
$H_0\cap V=\gen a$. Similarly $H_2\cap V=\gen c$ and
$H_3\cap V=\gen{ac}$. The objects of each nonedge type incident
with the edge $V$ are its two cosets modulo the indicated subgroup.
Any two such objects of different types meet inside $V$: two affine
lines of different directions in $C_2^2$ meet. Thus the residue of
$V$ is a complete three-partite graph with two objects of each type.
It has eight chambers, and every one of its flags extends. Translation
by $V$ acts freely on these eight choices (the three line kernels have
trivial intersection), with two orbits. Since $G$ is edge-transitive,
these give exactly two free global orbits. In the finite case there
are $8(|G|/4)=2|G|$ chambers.

Every pair of incident nonedge cosets has an actual common group
element $g$, so the edge $gV$ extends that pair. Singletons also
extend. The only remaining possible obstruction is therefore a
three-type flag avoiding edges, exactly as stated.
\end{proof}

The last condition is a genuine global incidence condition. It does
not follow merely from the description of a local Petrie turn.
If the starting rank-three truncation is the faithful thin incidence
map $M$, it reconstructs $M$ literally: deleting type 3 preserves
all its objects and incidences. In a quotient model with repeated
corners forgotten, this must instead be checked; it is not supplied
by group generation alone.

\begin{theorem}\label{classIII}
Let $M$ be a reflexible Class III map whose rank-three coset
truncation is its faithful thin incidence map. Then the four-type
pregeometry has trialities and no type-transposing correlations.
Its type-permutation quotient is exactly $C_3$. Its full-type chambers
have the two orbits of Lemma~\ref{local}. It is an incidence geometry
only if the additional completion condition in that lemma holds.
\end{theorem}
\begin{proof}
A marked triality cycles $a,c,ac$ and fixes $b$, so the three
periods are equal, say $m$, and the three nonedge parabolics have
order $2m$. None can be $G$: otherwise they all equal the same
dihedral group. Infinite dihedral groups have no subgroup $V$ of
order four. If its rotation subgroup has even order and $m>2$,
its central involution is characteristic and lies in $V$, contradicting
its cyclic permutation; if its rotation order is odd it has no $V$.
For $m=2$ the generators commute, and the existence of the marked
triality forces $b\notin V$. Then $G=C_2^3$, whose full $S_3$ on
$V\setminus1$ fixes $b$, giving Class IV. The case $m=1$ makes
$b=a$ and contradicts its being fixed while $a$ is cycled.
Thus $m>2$, and Lemma~\ref{local} applies.

An edge is incident with exactly two objects of each other type.
An object of any nonedge type is incident with $m$ edges, since
$[H_i:H_i\cap V]=m$ (an infinite cardinal if the period is infinite).
Since $m>2$, the edge type is the unique type whose objects have
exactly two neighbors of every other type. A correlation therefore
fixes the edge type, without needing finite cardinalities. The marked
triality induces an actual order-three coset permutation cycling the
other three types. If any correlation induced a transposition on
those types, conjugating by this triality would give one which swaps
faces and vertices while fixing the Petrie type. Its restriction to
the faithful rank-three truncation would be a duality of $M$.
By flag transitivity, compose it with a map automorphism so that it
fixes a chosen flag as a set while swapping its endpoint types.
Unique flag adjacencies then give a marked automorphism interchanging
$a,c$ and fixing $b$. This is a Wilson duality, contrary to Class III.
The image on types is consequently the displayed $C_3$. This argument
uses neither flag completion in rank four nor an unsupported equality
$\Aut(\Gamma)=G$.
\end{proof}

\section{A faithful toroidal obstruction to flag completion}

Take the square tessellation modulo the translations $(2,0),(0,2)$.
Write its vertices $v_{ij}$, horizontal edges $h_{ij}$ from $(i,j)$
to $(i+1,j)$, vertical edges $e_{ij}$ from $(i,j)$ to $(i,j+1)$,
and faces $f_{ij}$, with all indices modulo two. Parallel edges with
the same endpoints remain distinct. A face has boundary
\[
 h_{ij},\ e_{i+1,j},\ h_{i,j+1},\ e_{ij}.
\]
This is a two-cell torus map with four vertices, eight edges and
four square faces. Every incident vertex--face pair has two boundary
edges; every edge has two endpoints and two faces. Its 32 flags are
therefore faithful incidence triples, with unique flag adjacencies.

Let $D_8$ be the group of eight integral signed permutation matrices
in two coordinates and put $T=\F_2^2\rtimes D_8$. Crucially, we
retain the eight matrices as distinct: their action on the four
vertices alone is not faithful. Multiplication is
$(t,A)(u,B)=(t+Au,AB)$, with translations reduced modulo two.
Use
\[
 a=( (1,0),\operatorname{diag}(-1,1)),\quad
 b=(0,\begin{smallmatrix}0&1\\1&0\end{smallmatrix}),\quad
 c=(0,\operatorname{diag}(1,-1)).
\]
These are exactly the three reflections at the flag
$(v_{00},h_{00},f_{00})$. They generate $T$: $b,c$ generate
$D_8$, $abcb$ is translation $(1,0)$, and conjugating that
translation by $b$ gives $(0,1)$. Their action on the 32 flags is
regular. Conversely an incidence automorphism fixing one flag fixes
all flags by unique adjacencies and connectedness, so the full
incidence automorphism group is exactly $T$.

\begin{proposition}\label{torus}
The four-type coset pregeometry of this reflexible map is not a
geometry. In fact $(f_{00},v_{11},P)$ is an unextendable flag, where
\[
                 P=(h_{00},e_{11},h_{11},e_{00})
\]
is a Petrie circuit.
\end{proposition}
\begin{proof}
The displayed circuit alternately turns through the two sides of the
square tessellation, so is a Petrie polygon. It contains all four
vertices. The face $f_{00}$ and $P$ have precisely the edges
$h_{00},e_{00}$ in common, meeting at $v_{00}$. Hence they are
incident under both the literal coset definition and the source's
two-edge description. The opposite vertex $v_{11}$ belongs to
both polygons but to neither of those two common edges. Thus the
three objects are pairwise incident, while no edge completes them.
This is an actual surface-map counterexample to source Lemma 7,
not merely a degenerate quotient of its abstract map group.
\end{proof}

We will also need the Wilson symmetries and a precise word version
of this witness. Set $t=((1,1),I)$ and $j=(0,-I)$ in $T$. Direct
multiplication gives
\begin{equation}\label{words}
 (ab)^2=tj,\qquad (bc)^2=j,\qquad (acb)^2=t.
\end{equation}
The first two squares are independent central involutions. The
presentation by $a,b,c$ as involutions, $[a,c]=1$, with
$[a,b]$ and $[b,c]$ central, has at most 32 elements: commute
letters to the form $a^A c^C b^B[a,b]^X[b,c]^Y$, with five binary
exponents: central commutators of two involutions have order at most
two. Its actual image $T$ has 32 elements, so this is an
exact presentation. Every invertible linear change in $\gen{a,c}$
fixing $b$ preserves it. In particular
\[
 a\mapsto c,\quad c\mapsto ac,\quad b\mapsto b
\]
is a triality, and swapping $a,c$ is a duality. This torus map is
Class IV; it is not by itself a counterexample involving Class III.

For the coset version of Proposition~\ref{torus}, take
$F=H_0$, $Q=tH_2$, $P=H_3$. Here $t\in H_3$ and
$t=(ab)^2(bc)^2\in H_0H_2$, so all three pairs meet.
The geometric argument proves that no coset of $V$ meets all three.

\section{The obstruction persists for an actual Class III map}

\begin{theorem}\label{parallel}
Let $S$ be a finite perfect group with no nontrivial normal
2-subgroup, supporting a faithful regular Class III map with
markers $a_0,b_0,c_0$ and common odd period $m>2$.
Then the parallel product with the torus above is a faithful regular
Class III map. Its four-type coset pregeometry is not a geometry.
\end{theorem}
\begin{proof}
In $S\times T$ use the diagonal markers
$A=(a_0,a)$, $B=(b_0,b)$, $C=(c_0,c)$. They generate a subdirect
product. A proper subdirect product would give a nontrivial common
quotient of $S$ and $T$: project the two projection kernels and
identify the quotient using any lift. Such a quotient would be both
perfect and a finite 2-group, which is impossible. Thus the marked
group is $G=S\times T$.

For each nonedge parabolic the rotation in $S$ has order $m$ and
that in $T$ has order four. Since these orders are coprime, its cyclic
rotation group is their full direct product. The entire parabolic is
therefore precisely the pairs of dihedral elements with equal
rotation/reflection parity, and has order $8m$. The two original
end parabolics in each factor intersect in its marked $b$ subgroup.
Their fiber products consequently intersect in $\gen B$.
They each intersect $\gen{A,C}$ in the appropriate marked endpoint.
This is the rank-three intersection property: the resulting coset
map is thin with faithful flags. More explicitly, each edge has
two faces and two vertices, and its four choices are a single free
$\gen{A,C}$-orbit. An incident face--vertex pair has exactly two
edges because the end-parabolic intersection has order two; every
such pair extends by choosing an actual common group element.
Thus the flags are exactly the regular $G$-orbit, and their connected
adjacency graph is the Cayley graph of the generators. Gluing its flag
triangles yields a closed surface map, with dihedral face and vertex
residues. The same free-flag argument as above identifies its full
incidence automorphism group with $G$.

The two marked trialities act componentwise. If the product map had
a Wilson duality, the corresponding marked automorphism of $G$
would descend through the characteristic subgroup $O_2(G)=T$ to
a Wilson duality of the original $S$ map. The equality follows from
the hypothesis $O_2(S)=1$. This is impossible, so the product has
Class III. Its common period is $4m$.

Let $K_i$ denote its four parabolics, and put $z=(1,t)$.
Since $m$ is odd, \eqref{words} and the order-$m$ rotations give
\[
 (ACB)^{2m}=z,\qquad
 (AB)^{2m}=(1,tj),\qquad (BC)^{2m}=(1,j).
\]
Thus $z\in K_3\cap K_0K_2$ and
$(K_0,zK_2,K_3)$ is a pairwise incident flag. Projection to $T$
maps each coset of $K_i$ onto the corresponding coset of $H_i$
and preserves incidence. An edge completing this flag would
therefore complete the impossible torus flag $(H_0,tH_2,H_3)$.
There is none.
\end{proof}

An explicit instance is available without a large enumeration. Work in
$\F_8=\F_2[u]/(u^3+u+1)$ and use matrices
\[
 a_0=\begin{pmatrix}1&u\\0&1\end{pmatrix},\quad
 b_0=\begin{pmatrix}0&1\\1&0\end{pmatrix},\quad
 c_0=\begin{pmatrix}1&u^2\\0&1\end{pmatrix}.
\]
The complete small certificate accompanying this paper constructs all
504 determinant-one matrices, verifies that these three generate all
of them, checks the three periods nine and the end intersection
$\gen{b_0}$, and checks the field-squaring triality. Its finite
execution is a declared premise for this fixed instance, not an
exhaustion of maps. To exclude a marked duality it evaluates the
literal word
\[
 w=abababcbcbcbababc.
\]
The word is the identity on $(a_0,b_0,c_0)$; interchanging $a_0,c_0$
gives $\left(\begin{smallmatrix}u^2+u&u^2+u\\u&1\end{smallmatrix}\right)$,
which is not the identity. Triality then excludes the other two
Wilson transpositions as well.

For completeness the other group hypotheses have short proofs.
The group $\operatorname{SL}_2(8)$ is perfect: for a diagonal
matrix with diagonal entries $d,d^{-1}$, $d\ne1$, commutation
with upper unipotents gives every upper unipotent since $d^2-1\ne0$;
lower unipotents are conjugate, and elementary elimination shows these
generate the whole determinant-one group. Every involution is a
nonidentity unipotent and these form one conjugacy class. Its
centralizer has order eight, as direct two-by-two commutation shows,
so the class has 63 elements. A nontrivial normal 2-subgroup would
contain an involution and hence all 63, although a 2-subgroup of
this order-504 group has order at most eight. Thus $O_2(S)=1$.

The product example has group order $504\cdot32=16128$, common
period 36, 224 objects of each nonedge type and 4032 edges. It is
Class III and faithful, but its four-type structure is not a geometry.
Theorem~\ref{classIII} still proves that this pregeometry has a
triality and no type transposition. We do not misreport the failure
of flag completion as a failure of that separate symmetry conclusion.

\section{A separate failure of the unrestricted chamber-orbit claim}

For an odd $n\geq3$, take a dihedron: an $n$-cycle embedded on the
sphere with its two complementary $n$-gonal faces. This is a faithful
thin regular map with $4n$ flags and group $D_{2n}\times C_2$.
Choose adjacent reflecting generators $a,b$ for $D_{2n}$ and let
$c$ be the central involution swapping the faces. Then $ab$ has
order $n$, while $bc$ has order two and $(ac)b=(ab)c$ has order
$2n$. Consequently $H_3=\gen{b,ac}=G$: there is only one Petrie
polygon. The four-type structure is the original rank-three map
with a universal singleton type added. It is a geometry, and $G$
has exactly one chamber orbit, not two. This refutes source Lemma 8
in its unrestricted stated scope. The proper-parabolic hypothesis
of Lemma~\ref{local} is essential. No Class III exception is claimed.

These dihedra also correct two of the source's class correspondences.
For odd $n\geq3$, their periods $n,2,2n$ are distinct, so they are
Class I. In the four-type structure the vertex--edge incidence graph
is a $2n$-cycle; both faces and the singleton Petrie object are
incident with every object of the other types. Rotation by one step
of this incidence cycle interchanges vertices and edges, and fixes
the remaining objects. The type-permutation quotient is exactly
$C_2$, since the other type sizes are the distinct numbers two and
one. It is not the source's trivial Class I quotient.

For even $n\geq4$, the same dihedron has periods $n,2,n$. It is
Class II: the assignment $a\mapsto ac$, $b\mapsto b$, $c\mapsto c$
is an automorphism of $D_{2n}\times C_2$ because $n$ is even, while
the unequal period excludes the other permutations. Now $H_3$ has
index two and is different from $H_0$; both are normal. Also
$H_3\cap H_2=\gen b$. Consequently the two Petrie objects are
incident with all vertices, edges and faces. The four-type structure
has the same vertex--edge incidence cycle and two independent
universal pairs of face and Petrie objects. It admits the independent
type swaps vertex--edge and face--Petrie. Since $n>2$ distinguishes
these two pairs of type sizes, the exact type quotient is
$C_2\times C_2$, not the claimed Class II $C_2$.

\paragraph{Precise source disposition.}
Construction 6 defines a pregeometry for every marked group, but the
source's assertion that it always is a geometry is false, even for
a faithful Class III map. Lemma 7 is refuted by the torus and its
Class III parallel product. Lemma 8 is false without properness of
all three nonedge parabolics; its two full-type chamber orbits are
proved in Lemma~\ref{local} under that exact condition. The source's
universal path-residue explanation is unnecessary and false in small
maps. Root's separately reviewed digonal $C_2^3$ example addresses
Remark 9 and the claimed Class IV correlation quotient; that written
companion must be acquired in full if used in a complete L023 edition.
The odd and even dihedra refute the unrestricted Class I and II
quotient assertions, respectively.
Theorem 10's symmetry conclusion has the rigorous faithful-map
replacement in Theorem~\ref{classIII}; calling its output a geometry
requires the extra completion condition. Arbitrary quotient maps with
forgotten incidence flags are not silently treated as faithful maps.
There is no new conjecture, Lean declaration or numerical factor.
\end{document}
